% !TEX root = ../main.tex
\section{Introduction}
\label{sec:introduction}

\IEEEPARstart{W}{ildfires} threaten ecosystems, infrastructure, and human safety
\cite{bowman2020wildfire}. UAVs can rapidly observe hazardous and inaccessible
areas through flexible aerial sensing \cite{allan2022remote}, but operational
response requires observations to be converted into task zones, assigned across
a resource-limited fleet, and revised as the fire or UAV state changes.
Multi-UAV wildfire rescue is therefore a closed-loop decision problem rather
than an isolated perception or path-planning task.

An initial mission must be feasible for the current fleet: task priorities must
be matched to vehicle locations, endurance, payload, communication, and safety
limits \cite{floreano2015science}. Feasibility is also temporary because fire
evolution, resource consumption, communication changes, and UAV faults can
invalidate existing assignments. Initial resource-constrained mission planning
and runtime replanning must consequently be treated as coupled parts of the
same decision problem.

Existing work has advanced the individual capabilities needed for this process
\cite{erdelj2017help}. Thermal--RGB fusion, aerial semantic segmentation, and
learning-based 3D reconstruction improve fire detection and scene modeling
\cite{yuan2021uavfire,zhang2022uavtransformer,
yang2022uav3dreconstruction}, while autonomous exploration and mapping and
next-best-view planning update environmental knowledge and reduce observation
uncertainty \cite{popovic2021uavdisaster,bircher2016nbv}. On the action side,
deep reinforcement learning and coverage path planning support navigation and
large-area search \cite{xiao2021uavrl,yu2022coverageuav}, whereas multi-agent
reinforcement learning and human--robot collaboration support task allocation
and search-and-rescue operations
\cite{yang2023multiagentuav,waharte2021uavsar}. However, these capabilities are
commonly optimized and evaluated separately, often from a fixed initial state.
They therefore do not establish how visual and spatial evidence can be converted
into fire points and task zones for resource-constrained mission planning, or
how an initially feasible allocation should be revised when the fire, fleet, or
resource state changes.

Large language models (LLMs) provide a flexible interface between high-level
intent, structured observations, and robot behavior. SayCan grounds instructions
in executable skills \cite{ahn2022saycan}, Code as Policies generates
programmatic control policies \cite{liang2023codeaspolicies}, and related design
guidance exposes tools and environmental constraints to ChatGPT-generated robot
workflows \cite{vemprala2023chatgptrobotics}. Embodied and multimodal models
further connect reasoning to physical state: RT-1 learns vision-language-action
policies \cite{brohan2023rt1}, PaLM-E incorporates visual inputs and robot states
\cite{driess2023palme}, and Socratic Models compose specialized vision and
language models \cite{zeng2023socratic}. These approaches support task
decomposition and action sequencing, but they also show that generated actions
require grounding and external feasibility checks. Embodied and multimodal
models are generally evaluated on individual robots and predefined manipulation
or navigation tasks. Moving to a fleet-level decision layer introduces coupling
among vehicle states, task priorities, routes, and shared resources; linguistic
coherence therefore does not ensure that assignments jointly respect
endurance, payload, communication, priority, and safety constraints.

LLM-based agents have also used environmental feedback and accumulated
experience for sequential decision making. Inner Monologue revises embodied
plans from language-based feedback \cite{huang2023innermonologue}, LM-Nav uses
semantic waypoints for open-world navigation \cite{shah2023lmnav}, and Voyager
combines an LLM agent with an automatic curriculum and a growing skill library
\cite{wang2023voyager}. These systems primarily target single-agent embodied
environments rather than dynamic, interdependent tasks across a
resource-constrained fleet. They establish that feedback and accumulated
experience can improve sequential behavior, but do not establish how these
capabilities should be combined with fleet-wide resource constraints or how
their importance changes as operational pressure increases. Existing research
therefore provides the main ingredients for perception, planning, coordination,
and feedback, but not a common high-level decision loop that connects
fire-situation understanding, resource-constrained mission planning,
Validation, Memory, and Runtime Replanning. Strong perception does not by
itself establish that detected fire conditions can be converted into a feasible
fleet mission, while a feasible initial mission plan does not establish that the
mission can remain effective under subsequent perturbations. The central
question is whether LLMs can serve as the high-level closed-loop decision core
of a multi-UAV wildfire rescue system, and how Validation, Memory, and Runtime
Replanning sustain that role as operational pressure increases.

We address this question with an LLM-centered closed-loop decision framework
that transforms aerial observations through Vision Analysis and Mission
Planning into an initial fleet mission. Semantic Validation and Deterministic
Rule Validation reject plans that are tactically implausible or structurally
invalid, converting validation failures into targeted feedback for Mission
Planning. Episodic Memory retrieves concrete prior cases, whereas Lesson Memory
provides distilled guidance about recurring planning and resource-management
patterns. Runtime Replanning repairs the affected part of the active mission
after changes in the fire, fleet, resources, or schedule. Together, these
mechanisms turn one-shot LLM planning into a decision process that connects
fire-situation understanding, resource-constrained mission planning,
Validation, execution feedback, and Runtime Replanning within one continuous
loop.

We evaluate the framework in a real-UAV-informed AirSim suite through
end-to-end comparisons, controlled ablations, and compact-planner fine-tuning.

The main contributions of this work are summarized as follows:
\begin{enumerate}[leftmargin=1.6em,itemsep=2pt,topsep=3pt]
    \item We formulate resource-constrained multi-UAV wildfire rescue as a
    high-level closed-loop decision problem and develop an LLM-centered
    closed-loop decision framework that connects fire-situation understanding
    and resource-constrained mission planning with Validation, Memory, and
    Runtime Replanning.
    \item We construct a real-UAV-informed AirSim evaluation suite that exposes
    the decision loop to controlled increases in operational pressure through
    resource limitations, scheduling constraints, and runtime perturbations
    across four difficulty levels, and measures both task completion and
    execution quality with six mission-level metrics.
    \item We evaluate end-to-end behavior against an Optimization-Based Planner
    and across multiple foundation-model configurations, characterize through
    controlled ablations how Validation, Memory, and Runtime Replanning
    contribute as operational pressure increases, and assess whether validated
    demonstrations can specialize a compact planner with lower planning
    overhead.
\end{enumerate}

\endinput
